\documentclass[course]{course-report}

% 封面信息 - 请根据实际情况修改
\school{学校名称}{University Name}
\course{课程名称}
\title{课程报告题目}{Course Report Title}
\author{学生姓名}{Student Name}
\studentnumber{202600000000}
\advisor{指导老师姓名}{Advisor Name}
\department{学院名称}{Department Name}
\major{专业名称}{Major Name}

% 如果需要显示学校logo，请取消下面一行的注释
% 并将logo文件放在 pic/logo.pdf
% \showlogo

\begin{document}

\makecover

\thesistableofcontents

\chapter{引言}

\section{研究背景}

这里写课程报告的背景、问题来源和任务要求。课程报告通常不需要中英文摘要、原创性声明、插图目录、表格目录、符号表、致谢和成果清单。

\section{报告结构}

这里简要说明全文安排。本文后续章节演示公式、图表、交叉引用和参考文献的基本写法。

\chapter{方法与实现}

\section{基本方法}

正文可以直接使用中文、英文和数学公式。设目标函数为
\begin{equation}
  \label{eq:objective}
  J(\bm{x})=\sum_{i=1}^{n}(y_i-\bm{w}^{\mathrm{T}}\bm{x}_i)^2+\lambda\lVert\bm{w}\rVert_2^2 .
\end{equation}
式\ref{eq:objective}展示了一个带正则项的最小二乘目标。

\section{图表示例}

图\ref{fig:example}和表\ref{tab:example}展示了报告中常用的浮动体写法。

\begin{figure}[h]
  \centering
  \fbox{\parbox{0.6\textwidth}{\centering\vspace{2cm}这里是图片占位符\vspace{2cm}}}
  \caption{示例图片}
  \label{fig:example}
\end{figure}

\begin{table}[h]
  \caption{实验参数示例}
  \label{tab:example}
  \begin{tabular}{ccc}
    \toprule
    参数 & 取值 & 说明 \\
    \midrule
    $n$ & 100 & 样本数量 \\
    $\lambda$ & 0.01 & 正则化系数 \\
    迭代次数 & 50 & 最大迭代轮数 \\
    \bottomrule
  \end{tabular}
\end{table}

\chapter{结果与分析}

\section{实验结果}

这里描述实验设置、数据来源和主要结果。引用文献时可以继续使用模板提供的命令，例如\citing{example2024}。

\section{讨论}

这里分析结果的合理性、局限性和可能的改进方向。

\chapter{总结}

这里总结课程报告的主要工作和结论，并说明后续可以继续完善的方向。

% Uncomment to list all entries in reference.bib.
% \nocite{*}

\coursebibliography{reference}

\end{document}
